\documentclass[12pt]{article}

\usepackage[utf8]{inputenc}
\usepackage{amsmath, amssymb, amsthm}
\usepackage{geometry}
\geometry{a4paper, margin=2.5cm}

\title{Teorema Hahn--Banach}
\author{}
\date{}

\theoremstyle{definition}
\newtheorem{definition}{Definiție}

\theoremstyle{plain}
\newtheorem{proposition}{Propoziție}
\newtheorem{theorem}{Teoremă}

\begin{document}

\maketitle

\begin{definition}
Fie $X$ un spațiu vectorial real și $p : X \to \mathbb{R}$ o funcție subliniară, adică


\[
p(x+y) \le p(x) + p(y), \qquad p(\lambda x) = \lambda p(x) \quad \text{pentru } \lambda \ge 0.
\]


Dacă $Y \subset X$ este un subspațiu și $f_0 : Y \to \mathbb{R}$ este o aplicație liniară, spunem că $f_0$ este \emph{dominată} de $p$ dacă


\[
f_0(y) \le p(y) \quad \text{pentru toate } y \in Y.
\]


\end{definition}

\begin{theorem}[Hahn--Banach, forma reală]
Fie $X$ un spațiu vectorial real, $p : X \to \mathbb{R}$ o funcție subliniară și $Y \subset X$ un subspațiu. Dacă $f_0 : Y \to \mathbb{R}$ este liniară și dominată de $p$, atunci există o aplicație liniară $f : X \to \mathbb{R}$ cu proprietățile:
\begin{enumerate}
    \item[(i)] $f$ extinde pe $f_0$, adică $f|_Y = f_0$;
    \item[(ii)] $f$ este dominată de $p$, adică
    

\[
    f(x) \le p(x) \quad \text{pentru toate } x \in X.
    \]


\end{enumerate}
\end{theorem}

\begin{proposition}
În ipotezele teoremei Hahn--Banach, următoarele afirmații sunt echivalente pentru o extensie $f : X \to \mathbb{R}$:
\begin{enumerate}
    \item[(i)] $f$ este o extensie liniară a lui $f_0$ dominată de $p$.
    \item[(ii)] Există o funcție $\psi : X \to \mathbb{R}$, cu $\psi(x) \ge 0$ pentru toate $x$, astfel încât
    

\[
    f(x) = f_0(y) + \psi(x-y)
    \]


    pentru orice descompunere $x = y + (x-y)$ cu $y \in Y$, iar $\psi$ satisface
    

\[
    \psi(z) \le p(z) \quad \text{pentru toate } z \in X.
    \]


\end{enumerate}
\end{proposition}

\end{document}
